% Results for the J-BHI paper. Every number is a macro from generated/c2c1_numbers.tex or
% generated/c2c2_numbers.tex (scripts/paper_c2c1.py, scripts/paper_c2c2.py; inputs in results/).

\section{Results}\label{sec:results}

\subsection{Pipeline Validation}\label{sec:res_pipeline}
Table~\ref{tab:gates} lists every gate of both registrations. Our single-cell processing
reproduced the published per-line TH$^+$ fractions almost exactly ($r = \cconeReproR$,
$n = \cconeReproN$ lines; A0-6 pass). The positive control was harder to meet. A first version of
A0-4, which predicted our own TH$^+$ fraction, failed (CV AUC \cconeAucv). An exact replication of
the published classifier (\ref{eq:enet}) on \cconeDayZeroGenes{} expressed genes reached a mean CV
AUC of \cconeAuc{} (SD \cconeAucSD; permutation $p = \cconeAucP$; \cconeHigh{} high- and
\cconeLow{} low-yield lines), against 0.85 reported in \cite{bressan2023foundin} and
\cconeAucBatch{} for a batch-only model. This is the ``partial'' verdict (AUC between 0.65 and
0.75 with $p < 0.05$). The analysis proceeded, with the shortfall reported as a limitation.

% AUTO-GENERATED by scripts/paper_c2c1.py from results/ -- do not edit by hand.
\begin{table*}[!t]
\centering
\caption{Gates of both registrations: rule fixed before the data were seen, observed value, verdict.}
\label{tab:gates}
\footnotesize
\begin{tabular}{p{2.9cm}p{5.4cm}p{5.6cm}p{2.4cm}}
\toprule
Gate & Rule & Observed & Verdict \\
\midrule
\multicolumn{4}{l}{\emph{Registration 1 (all-cell phenotypes)}} \\
A0-6 Reproduction & $r(\%\mathrm{TH}^+) \ge 0.80$ vs.\ published & $r = 0.999$ ($n=79$) & pass \\
A0-4 Positive control & day-0 CV AUC $\ge 0.75$, $p<0.05$ & AUC 0.72, $p=0.020$ & partial \\
A0-2 Stability & $\mathrm{ICC} \ge 0.5$ & P1 0.33, P2 0.94, P3 0.55, P4 0.64 & P1 fail; P2--P4 pass \\
A0-3 Cells & $\ge 1{,}000$ QC cells per run & 1 line excluded & pass \\
A0-1 Linkage & $\ge 30$ donors with $\ge 3$ OFF visits over $\ge 2$\,y & 44/47 & pass \\
A0-5 Coherence & $\ge 80\%$ member genes share sign & P3 75\%, P4 62\% & fail (both) \\
\multicolumn{4}{l}{\emph{Registration II (DA-neuron phenotypes)}} \\
G-cells & $\ge 100$ DA cells per run & 0 lines excluded & pass \\
G-icc & FOUNDIN $\ge 0.5$ and Jerber $\ge 0.40$ & D1 $-$0.24/0.38, D2 0.83/0.62, D3 0.46/0.18, D4 0.88/0.43 & D2, D4 pass \\
G-eQTL & concordance $\ge 0.70$, $\pi_1 \ge 0.30$, both $>$ null P95 & 0.96, 0.79 ($m=28$) & pass \\
A0-1 Linkage & as registration 1 & 44/47 & pass \\
A0-5 Coherence (D4) & $\ge 80\%$ member genes share sign & 59\% of 78 & fail \\
\bottomrule
\end{tabular}
\end{table*}


\subsection{Registration 1: All-Cell Phenotypes}\label{sec:res_reg1}
\emph{Gates.} Technical stability (\ref{eq:icc}) removed the DA-fraction phenotype P1
($\widehat{\mathrm{ICC}} = \cconeIccPone$) and retained SNCA, OXPHOS and lysosome (P2--P4:
\cconeIccPtwo, \cconeIccPthree, \cconeIccPfour). A0-3 excluded \cconeAthreeExcl{} donor line,
leaving \cconeN{} PD donors with linked PPMI data, of whom \cconeAzeroOne{} met the clinical rule
(\ref{eq:a01}). Both pathway scores then failed coherence (\ref{eq:a05}): only
\cconeAfivePthree\% of \cconeAfivePthreeN{} OXPHOS genes and \cconeAfivePfour\% of
\cconeAfivePfourN{} lysosome genes shared the composite's sign, so P3 and P4 were not tested on
UPDRS~III.

\emph{Tests.} None of the \cconeNtests{} confirmatory tests was supported
(Fig.~\ref{fig:reg1}): \cconeNeq{} were equivalent to zero within $\pm 0.40$ and \cconeNinc{} were
inconclusive. For the motor outcome, all-cell SNCA gave $\hat\beta_3 = \cconeHoneBeta$
($p = \cconeHoneP$; partial $r = \cconeHoneR$, 90\% CI \cconeHoneRlo{} to \cconeHoneRhi;
\cconeVisits{} visits). The OXPHOS estimate removed by A0-5 was larger ($r = \cconeHthreeR$,
$p = \cconeHthreeP$) and is reported descriptively only. All three SBR tests
($n = \cconeSBRn$ donors, \cconeSBRobs{} scans) were equivalent. No registered sensitivity analysis
produced a supported result (\cconeSensSup{} of all sensitivity tests).

\begin{figure}[!t]
\centering
\includegraphics[width=\columnwidth]{fig3_forest.pdf}
\caption{Registration~1: partial correlation $r$ (\ref{eq:partialr}) with 90\% CI for the
confirmatory tests of the all-cell phenotypes. Shaded: equivalence region $\pm 0.40$
(\ref{eq:tost}). Class (\ref{eq:class}) is shown by colour and marker shape.}
\label{fig:reg1}
\end{figure}

\subsection{DA-Neuron Labels and Registration-II Gates}\label{sec:res_reg2gates}
\emph{Labels.} Among neurons, the old rule (TH $\ge 1$ UMI) was not DA-specific: apart from TH,
canonical DA markers were not enriched relative to other neurons. Under the new rule
(\ref{eq:dalabel}) DA neurons were enriched for every DA marker (median log$_2$ ratio against other
neurons: TH \cctwoTH; SLC6A3, the weakest, \cctwoSLCthree) and depleted of the GABAergic marker
GAD1 (\cctwoGADone). The calibration (\ref{eq:calib}) selected $t^\star = \cctwoT$, the lower edge
of the grid $[\cctwoGridMin, \cctwoGridMax]$. Agreement with the published DA fractions was lower
than for the old rule ($r = \cctwoCalR$ vs.\ \cctwoCalRold), and kept rising as the rule broadened.
The published fractions therefore track how many neurons a run makes rather than DA identity, and
we chose the new label for specificity, not for agreement. The median run had \cctwoDAcells{} DA
neurons, and \cctwoDAlow{} runs fell below the G-cells threshold.

\emph{Stability.} FOUNDIN control-run ICCs and the across-batch ICCs from the independent dataset
were \cctwoIccDone/\cctwoJerDone{} for D1, \cctwoIccDtwo/\cctwoJerDtwo{} for D2 (across-batch
95\% CI \cctwoJerDtwoLo{} to \cctwoJerDtwoHi), \cctwoIccDthree/\cctwoJerDthree{} for D3 and
\cctwoIccDfour/\cctwoJerDfour{} for D4. D2 and D4 passed G-icc; D1 and D3 were not tested and are
reported descriptively.

\emph{Genetic positive control.} Of \cctwoLeads{} lead eQTLs, \cctwoExpressed{} genes were
expressed and \cctwoOnArray{} lead SNPs were directly genotyped, leaving $m = \cctwoTested$
testable pairs. In the new DA-neuron pseudobulk (\ref{eq:eqtl}), sign concordance was
$\kappa = \cctwoConc$ and $\hat\pi_1 = \cctwoPione$ (\ref{eq:pi1}), above the null 95th percentiles
(\cctwoNullConc{} and \cctwoNullPione). G-eQTL passed. The old-label ($\kappa = \cctwoOldConc$,
$\hat\pi_1 = \cctwoOldPione$, $m = \cctwoOldTested$) and all-cell ($\kappa = \cctwoAllConc$,
$\hat\pi_1 = \cctwoAllPione$, $m = \cctwoAllTested$) pseudobulks passed as well. The gate thus
shows that genetic regulation is present in the new label, not that the new label is more
DA-specific than the alternatives.

\emph{Clinical sufficiency and coherence.} All \cctwoN{} PD donors passed G-cells and
\cctwoAzeroOne{} met (\ref{eq:a01}). The lysosome score D4 failed coherence
(\cctwoAfiveShare\% of \cctwoAfiveGenes{} genes), so H1 on D4 was not tested
(descriptive $r = \cctwoDfourR$, $p = \cctwoDfourP$); the secondary family kept \cctwoSecN{} tests.

\subsection{Registration II: Confirmatory Tests}\label{sec:res_reg2}
\emph{Primary (H1-DA).} Higher SNCA expression in DA neurons was associated with slower motor
progression: $\hat\beta_3 = \cctwoBeta$ UPDRS~III points per year per SD (SE \cctwoSE,
$\nu = \cctwoDF$, $p = \cctwoP$), partial $r = \cctwoR$ (90\% CI \cctwoRlo{} to \cctwoRhi;
$n = \cctwoN$ donors, \cctwoVisits{} visits). Under (\ref{eq:class}) the primary hypothesis is
supported. Its CI excludes zero but does not exclude correlations near zero, and it overlaps the
equivalence region (Fig.~\ref{fig:reg2}).

\emph{Secondary and H4.} No secondary test was supported. The largest was SNCA by SAA status
($r = \cctwoSaaR$, $p = \cctwoSaaP$, $q = \cctwoSaaQ$; inconclusive). Bulk OXPHOS (H4) was
equivalent to zero ($r = \cctwoHfourR$, $p = \cctwoHfourP$), so the registration-1 exploratory OXPHOS
signal did not carry over to an independent assay of the same lines.

\emph{Sensitivity.} The direction of the primary effect held in every registered sensitivity
analysis, but its strength did not:
\begin{itemize}
\item old DA label: $\hat\beta_3 = \cctwoOldBeta$, $p = \cctwoOldP$;
\item composition-adjusted: $\hat\beta_3 = \cctwoAdjBeta$, $p = \cctwoAdjP$;
\item idiopathic PD only: $p = \cctwoIdioP$ ($n = \cctwoIdioN$);
\item five-level batch dummies: $\hat\beta_3 = \cctwoBatchBeta$, $p = \cctwoBatchP$;
\item visits within 5 years of diagnosis: $p = \cctwoFiveP$;
\item leaving out one donor gave $\hat\beta_3$ from \cctwoLooMin{} to \cctwoLooMax, with
\cctwoLooSign{} sign changes.
\end{itemize}
The batch-adjusted estimate is smaller and no longer significant.

\begin{figure*}[!t]
\centering
\includegraphics[width=0.85\textwidth]{fig5_c2c2_forest.pdf}
\caption{Registration~II: partial correlation $r$ (\ref{eq:partialr}) with 90\% CI for the primary
test, its registered sensitivity analyses, the secondary family, H4 and the phenotypes not tested.
Class (\ref{eq:class}) is shown by colour and marker shape; shaded: equivalence region
$\pm 0.40$.}
\label{fig:reg2}
\end{figure*}

\subsection{Exploratory, Post Hoc: Batch Structure}\label{sec:res_batch}
The decomposition (\ref{eq:within}) was specified after the primary result was seen and is
exploratory. D2 differed between batches ($R^2 = \cctwoBatchRtwoDtwo$, $p = \cctwoBatchPDtwo$);
per-donor progression slopes differed less ($R^2 = \cctwoBatchRtwoSlope$,
$p = \cctwoBatchPSlope$), so the pre-set confounding criterion was not met. The simple per-donor
association was $r = \cctwoRall$ overall and $r_W = \cctwoRwithin$ within batches
($p = \cctwoRwithinP$), which meets the pre-set rule for an association lying mainly within batches.
The five batch means pointed the same way ($r_B = \cctwoRbetween$), so a partial between-batch
contribution cannot be excluded. The registered medium covariate captured only
\cctwoMediaShareD\% (D2) and \cctwoMediaShareS\% (slopes) of the between-batch variance, which
explains why replacing it with batch dummies changed the estimate.
